%% kapa: exam=1 round=19 year=2008 period=1 form=A subject=appraisal_law subject_ko=부동산관계법규 seq=4 number=44
%% source: Q-Net 감정평가사 기출문제 (cst00309) artlSeq=1001962 file=extracted/1008/1001962/2008년 제19회 감정평가사 1차시험 기출문제/1교시.hwp
%% conversion: kordoc 4.24.51 -> tools/gen1.py (auto); review=proofread-sonnet; flags=none
\begin{question}{44}
국토의 계획 및 이용에 관한 법령상 도시기본계획에 관한 내용으로 옳은 것은?
\begin{choices}
  \item 도시기본계획수립권자는 특별시장․광역시장․시장 또는 군수이나 그 수립기준은 국토해양부장관이 정한다.
  \item 도시기본계획과 광역도시계획의 내용이 다를 경우 도시기본계획의 내용이 우선한다.
  \item 도시기본계획을 수립․변경하는 경우에는 미리 당해 지방의회의 승인을 받아야 한다.
  \item 도시기본계획수립권자는 10년마다 관할구역의 도시기본계획의 타당성 여부를 전반적으로 재검토하여 정비해야 한다.
  \item 도시기본계획을 수립 또는 변경하는 경우에는 공청회를 개최할 필요가 없다.
\end{choices}
\end{question}
\answer{1}
